\documentclass{amsart}
\usepackage[UTF8,fontset=windows]{ctex}
\usepackage[a4paper,margin=24mm]{geometry}
% For dealing with references we use the comment environment
\usepackage{verbatim}
\newenvironment{reference}{\comment}{\endcomment}
%\newenvironment{reference}{}{}
\newenvironment{slogan}{\comment}{\endcomment}
\newenvironment{history}{\comment}{\endcomment}

% For commutative diagrams we use Xy-pic
\usepackage[all]{xy}

% We use 2cell for 2-commutative diagrams.
\xyoption{2cell}
\UseAllTwocells

% We use multicol for the list of chapters between chapters
\usepackage{multicol}

% This is generally recommended for better output
\usepackage{lmodern}
\usepackage[T1]{fontenc}

% For cross-file-references

% Package for hypertext links:
\usepackage{hyperref}

% For any local file, say "hello.tex" you want to link to please

% Theorem environments.
%
\theoremstyle{plain}
\newtheorem{theorem}[subsection]{Theorem}
\newtheorem{proposition}[subsection]{Proposition}
\newtheorem{lemma}[subsection]{Lemma}

\theoremstyle{definition}
\newtheorem{definition}[subsection]{Definition}
\newtheorem{example}[subsection]{Example}
\newtheorem{exercise}[subsection]{Exercise}
\newtheorem{situation}[subsection]{Situation}

\theoremstyle{remark}
\newtheorem{remark}[subsection]{Remark}
\newtheorem{remarks}[subsection]{Remarks}

\numberwithin{equation}{subsection}

% Macros
%
\def\lim{\mathop{\mathrm{lim}}\nolimits}
\def\colim{\mathop{\mathrm{colim}}\nolimits}
\def\Spec{\mathop{\mathrm{Spec}}}
\def\Hom{\mathop{\mathrm{Hom}}\nolimits}
\def\Ext{\mathop{\mathrm{Ext}}\nolimits}
\def\SheafHom{\mathop{\mathcal{H}\!\mathit{om}}\nolimits}
\def\SheafExt{\mathop{\mathcal{E}\!\mathit{xt}}\nolimits}
\def\Sch{\mathit{Sch}}
\def\Mor{\mathop{\mathrm{Mor}}\nolimits}
\def\Ob{\mathop{\mathrm{Ob}}\nolimits}
\def\Sh{\mathop{\mathit{Sh}}\nolimits}
\def\NL{\mathop{N\!L}\nolimits}
\def\CH{\mathop{\mathrm{CH}}\nolimits}
\def\proetale{{pro\text{-}\acute{e}tale}}
\def\etale{{\acute{e}tale}}
\def\QCoh{\mathit{QCoh}}
\def\Ker{\mathop{\mathrm{Ker}}}
\def\Im{\mathop{\mathrm{Im}}}
\def\Coker{\mathop{\mathrm{Coker}}}
\def\Coim{\mathop{\mathrm{Coim}}}

% Boxtimes
%
\DeclareMathSymbol{\boxtimes}{\mathbin}{AMSa}{"02}

%
% Macros for moduli stacks/spaces
%
\def\QCohstack{\mathcal{QC}\!\mathit{oh}}
\def\Cohstack{\mathcal{C}\!\mathit{oh}}
\def\Spacesstack{\mathcal{S}\!\mathit{paces}}
\def\Quotfunctor{\mathrm{Quot}}
\def\Hilbfunctor{\mathrm{Hilb}}
\def\Curvesstack{\mathcal{C}\!\mathit{urves}}
\def\Polarizedstack{\mathcal{P}\!\mathit{olarized}}
\def\Complexesstack{\mathcal{C}\!\mathit{omplexes}}
% \Pic is the operator that assigns to X its picard group, usage \Pic(X)
% \Picardstack_{X/B} denotes the Picard stack of X over B
% \Picardfunctor_{X/B} denotes the Picard functor of X over B
\def\Pic{\mathop{\mathrm{Pic}}\nolimits}
\def\Picardstack{\mathcal{P}\!\mathit{ic}}
\def\Picardfunctor{\mathrm{Pic}}
\def\Deformationcategory{\mathcal{D}\!\mathit{ef}}

\hypersetup{hidelinks}
\widowpenalty=10000\clubpenalty=10000\displaywidowpenalty=10000
\begin{document}
\title[Stacks Project, Tag 0FLB]{有限 syntomic 态射的微分形式迹}
\maketitle
\noindent\textbf{Counting scope.} The primary problem is Tag 0FLB, marked below. All preceding numbered statements are supporting author context and are not counted as additional problems. This is a local pilot with a temporary ID, not production admission.


\noindent Copyright (C) 2005--2025 Johan de Jong. Stacks Project authors. Source: \href{https://stacks.math.columbia.edu/tag/0FLB}{Tag 0FLB}.

\noindent Difficulty H2: 本题构造与基变换相容、在拉回微分形式上限制为普通迹的微分形式迹映射。证明先用upstairs映射及Tate同构获得Noetherian情形，再在通用有限syntomic模型上证明唯一性并以Noetherian逼近扩展至任意基底。构造所用核心映射是先修输入，但唯一性、极限定义及基变换相容仍由本证明组织，判H2。.

\noindent History: excerpt from revision local cached snapshot; exact chapter hashes; revision unknown, derham.tex, lines 3820--3968. Only reference presentation was adapted. Licensed under GNU FDL 1.2 or later; no invariant sections or cover texts. See accompanying COPYING.

\begin{lemma}
\label{lemma-Garel-upstairs}
There exists a unique rule that to every locally quasi-finite syntomic
morphism of schemes $f : Y \to X$ assigns $\mathcal{O}_Y$-module maps
$$
c^p_{Y/X} :
\Omega^p_{Y/\mathbf{Z}}
\longrightarrow
f^*\Omega^p_{X/\mathbf{Z}} \otimes_{\mathcal{O}_Y} \det(\NL_{Y/X})
$$
for $p \geq 0$ satisfying (\href{https://stacks.math.columbia.edu/tag/0H9C}{item degree zero (Tag 0H9C)}), (\href{https://stacks.math.columbia.edu/tag/0H9D}{item multiplicative (Tag 0H9D)}),
and (\href{https://stacks.math.columbia.edu/tag/0H9E}{item base change (Tag 0H9E)}). In particular, the composition of
$c^p_{Y/X}$ with $f^*\Omega^p_{X/\mathbf{Z}} \to \Omega^p_{Y/\mathbf{Z}}$
is multiplication by $\delta(\NL_{Y/X})$.
\end{lemma}

\begin{proof}
This proof is very similar to the proof of
Discriminants, Proposition \ref{discriminant-proposition-tate-map}
and we suggest the reader look at that proof first.

\medskip\noindent
Let us reformulate the statement. Consider the category
$\mathcal{C}$ whose objects, denoted $Y/X$, are locally quasi-finite syntomic
morphism $f : Y \to X$ of schemes and whose morphisms
$b/a : Y'/X' \to Y/X$ are commutative diagrams
$$
\xymatrix{
Y' \ar[d]_{f'} \ar[r]_b & Y \ar[d]^f \\
X' \ar[r]^a & X
}
$$
which induce an isomorphism of $Y'$ with an open subscheme of
$X' \times_X Y$. The lemma means that for every object
$Y/X$ of $\mathcal{C}$ we have maps $c^p_{Y/X}$, $p \geq 0$
with properties (\href{https://stacks.math.columbia.edu/tag/0H9C}{item degree zero (Tag 0H9C)}) and (\href{https://stacks.math.columbia.edu/tag/0H9D}{item multiplicative (Tag 0H9D)}),
and for every morphism $b/a : Y'/X' \to Y/X$ of $\mathcal{C}$ we have
that $c^p_{Y/X}$ and $c^p_{Y'/X'}$ are compatible as in
Lemma \href{https://stacks.math.columbia.edu/tag/0H9G}{lemma base change Garel upstairs (Tag 0H9G)}.

\medskip\noindent
Given $Y/X$ in $\mathcal{C}$ and $y \in Y$ we can find
an affine open $V \subset Y$ and $U \subset X$ with $f(V) \subset U$
such that there exists some maps
$$
\Omega^p_{Y/\mathbf{Z}}|_V
\longrightarrow
\left(
f^*\Omega^p_{X/\mathbf{Z}} \otimes_{\mathcal{O}_Y} \det(\NL_{Y/X})
\right)|_V
$$
with properties (\href{https://stacks.math.columbia.edu/tag/0H9C}{item degree zero (Tag 0H9C)}) and (\href{https://stacks.math.columbia.edu/tag/0H9D}{item multiplicative (Tag 0H9D)}).
This follows from picking affine opens as in
Discriminants, Lemma \href{https://stacks.math.columbia.edu/tag/0BWE}{lemma syntomic quasi finite (Tag 0BWE)} part (5)
and then using the construction in Remark \href{https://stacks.math.columbia.edu/tag/0FL9}{remark local description (Tag 0FL9)}.

\medskip\noindent
Note that the \'etale locus of $f$ is exactly the set of points where
$\delta(\NL_{Y/X})$ does not vanish, see discussion in
Discriminants, Section \href{https://stacks.math.columbia.edu/tag/0FKB}{section tate map (Tag 0FKB)}.
In particular
$f^*\Omega^p_{X/\mathbf{Z}} \to \Omega^p_{Y/\mathbf{Z}}$
becomes an isomorphism after inverting $\delta(\NL_{Y/X})$.
We conclude that if $\Omega^p_{X/\mathbf{Z}}$ is finite locally free and
the annihilator of $\delta(\NL_{Y/X})$ in $\mathcal{O}_Y$ is zero, then
the local maps constructed in the previous paragraph
are unique and automatically glue!

\medskip\noindent
Let $\mathcal{C}_{nice} \subset \mathcal{C}$ denote the full subcategory
of $Y/X$ such that
\begin{enumerate}
\item $\Omega_{X/\mathbf{Z}}$ is locally free, and
\item the annihilator of $\delta(\NL_{Y/X})$ in $\mathcal{O}_Y$ is zero.
\end{enumerate}
By the remarks in the previous paragraph, we see that for any
object $Y/X$ of $\mathcal{C}_{nice}$ we have unique maps
$c^p_{Y/X}$, $p \geq 0$ satisfying conditions (\href{https://stacks.math.columbia.edu/tag/0H9C}{item degree zero (Tag 0H9C)})
and (\href{https://stacks.math.columbia.edu/tag/0H9D}{item multiplicative (Tag 0H9D)}). If $b/a : Y'/X' \to Y/X$
is a morphism of $\mathcal{C}_{nice}$, then the maps
$c^p_{Y/X}$ and $c^p_{Y'/X'}$ are compatible as in
Lemma \href{https://stacks.math.columbia.edu/tag/0H9G}{lemma base change Garel upstairs (Tag 0H9G)}: namely,
locally there do exist compatible maps by
Lemma \href{https://stacks.math.columbia.edu/tag/0H9H}{lemma existence base change Garel upstairs (Tag 0H9H)}
and the uniqueness just mentioned shows these maps agree
with $c^p_{Y/X}$ and $c^p_{Y'/X'}$.
In other words, we have solved the problem
on the full subcategory $\mathcal{C}_{nice}$. For $Y/X$ in $\mathcal{C}_{nice}$
we continue to denote $c^p_{Y/X}$ the solution we've just found.

\medskip\noindent
In fact, more generally, suppose we have a morphism $b/a : Y'/X' \to Y/X$
of $\mathcal{C}$ such that $Y/X$ is an object of $\mathcal{C}_{nice}$.
The same argument, but this time using the uniqueness in
Lemma \href{https://stacks.math.columbia.edu/tag/0H9G}{lemma base change Garel upstairs (Tag 0H9G)}, tells us there are
maps $c^p_{Y'/X'}$, $p \geq 0$  satisfying conditions (\href{https://stacks.math.columbia.edu/tag/0H9C}{item degree zero (Tag 0H9C)})
and (\href{https://stacks.math.columbia.edu/tag/0H9D}{item multiplicative (Tag 0H9D)}) compatible with the already constructed
maps $c^p_{Y/X}$. Let us call this the base change
$(b/a)^*c^p_{Y/X}$.

\medskip\noindent
Consider morphisms
$$
Y_1/X_1 \xleftarrow{b_1/a_1} Y/X \xrightarrow{b_2/a_2} Y_2/X_2
$$
in $\mathcal{C}$ such that $Y_1/X_1$ and $Y_2/X_2$ are objects
of $\mathcal{C}_{nice}$.
{\bf Claim.} The two base changes $(b_i/a_i)^*c^p_{Y_i/X_i}$, $i = 1, 2$
are equal. We will first show that the claim implies the lemma
and then we will prove the claim.

\medskip\noindent
Let $d, n \geq 1$ and consider the locally
quasi-finite syntomic morphism $Y_{n, d} \to X_{n, d}$
constructed in Discriminants, Example
\href{https://stacks.math.columbia.edu/tag/0FK8}{example universal quasi finite syntomic (Tag 0FK8)}.
Then $Y_{n, d}$ and $Y_{n, d}$ are irreducible schemes of finite type and
smooth over $\mathbf{Z}$. Namely, $X_{n, d}$ is a spectrum of a
polynomial ring over $\mathbf{Z}$ and $Y_{n, d}$ is an open subscheme
of such. The morphism $Y_{n, d} \to X_{n, d}$ is locally quasi-finite syntomic
and \'etale over a dense open, see Discriminants, Lemma
\href{https://stacks.math.columbia.edu/tag/0FK9}{lemma universal quasi finite syntomic etale (Tag 0FK9)}.
Thus $\delta(\NL_{Y_{n, d}/X_{n, d}})$ is nonzero: for example we have
the local description of $\delta(\NL_{Y/X})$ in
Discriminants, Remark \href{https://stacks.math.columbia.edu/tag/0FKC}{remark local description delta (Tag 0FKC)}
and we have the local description of \'etale morphisms in
Morphisms, Lemma \href{https://stacks.math.columbia.edu/tag/02GU}{lemma etale at point (Tag 02GU)} part (8).
Now a nonzero section of an invertible module over an irreducible
regular scheme has vanishing annihilator. Thus
$Y_{n, d}/X_{n, d}$ is an object of $\mathcal{C}_{nice}$.

\medskip\noindent
Let $Y/X$ be an arbitrary object of $\mathcal{C}$. Let $y \in Y$.
By Discriminants, Lemma \href{https://stacks.math.columbia.edu/tag/0FKA}{lemma locally comes from universal (Tag 0FKA)}
we can find $n, d \geq 1$ and morphisms
$$
Y/X \leftarrow V/U \xrightarrow{b/a} Y_{n, d}/X_{n, d}
$$
of $\mathcal{C}$ such that $V \subset Y$ and $U \subset X$ are open.
Then we have the base change $c^p_{V/U} = (b/a)^*c^p_{Y_{n, d}/X_{n, d}}$.
The claim guarantees these locally constructed maps $c^p_{V/U}$ glue!
Thus we get a well defined global maps $c^p_{Y/X}$ with properties
(\href{https://stacks.math.columbia.edu/tag/0H9C}{item degree zero (Tag 0H9C)}) and (\href{https://stacks.math.columbia.edu/tag/0H9D}{item multiplicative (Tag 0H9D)}).
Finally, let $b/a : Y'/X' \to Y/X$ be an arbitray morphism of $\mathcal{C}$.
We have the maps $c^p_{Y'/X'}$, $p \geq 0$ and $c^p_{Y/X}$, $p \geq 0$
constructed in this paragraph. To check they are compatible as in
Lemma \href{https://stacks.math.columbia.edu/tag/0H9G}{lemma base change Garel upstairs (Tag 0H9G)} we may work locally
on $Y'/X'$ and $Y/X$. Thus we may assume that there exists a morphism
$Y/X \to Y_{n, d}/X_{n, d}$. By construction, both the family
$c^p_{Y'/X'}$ and the family $c^p_{X/Y}$ are compatible with
the map towards $Y_{n, d}/X_{n, d}$. It follows easily from this
(and the uniqueness in Lemma \href{https://stacks.math.columbia.edu/tag/0H9G}{lemma base change Garel upstairs (Tag 0H9G)})
that $c^p_{Y'/X'}$ is compatible with $c^p_{Y/X}$.
Thus it remains to prove the claim.

\medskip\noindent
In the rest of the proof we prove the claim. We may pick a point
$y \in Y$ and prove the maps agree in an open neighbourhood of $y$.
Thus we may replace $Y_1$, $Y_2$ by open neighbourhoods of the
image of $y$ in $Y_1$ and $Y_2$. Therefore we may assume
$Y, X, Y_1, X_1, Y_2, X_2$ are affine, say they are the spectra
of rings $B, A, B_1, A_1, B_2, A_2$. Picture
$$
\xymatrix{
B_1 \ar[r] & B & B_2 \ar[l] \\
A_1 \ar[u] \ar[r] & A \ar[u] & A_2 \ar[l] \ar[u]
}
$$
By assumption the spectrum of $B$ is an affine open of
both the spectrum of $A \otimes_{A_1} B_1$ and $A \otimes_{A_2} B_2$.
Shrinking more we may assume there exist elements
$g_i \in A \otimes_{A_i} B_i$ such that our maps give isomorphisms
$(A \otimes_{A_i} B_i)_{g_i} = B$, see
Properties, Lemma \href{https://stacks.math.columbia.edu/tag/0H9B}{lemma standard open two affines (Tag 0H9B)}.
Let $x_\alpha$, $y_\beta$ be a sufficiently large collection of
variables such that we may choose surjections
$$
A'_1 = A_1[x_\alpha] \to A
\quad\text{and}\quad
A'_2 = A_2[y_\beta] \to A
$$
of $A_1$ and $A_2$-algebras. Then we can choose lifts
$h_i \in A'_i \otimes_{A_i} B_i$ of $g_i \in A \otimes_{A_i} B_i$
and we consider the diagram
$$
\xymatrix{
(A'_1 \otimes_{A_1} B_1)_{h_1} \ar[r] & B &
(A'_2 \otimes_{A_1} B_2)_{h_2} \ar[l] \\
A'_1 \ar[u] \ar[r] & A \ar[u] & A'_2 \ar[l] \ar[u]
}
$$
By construction the two squares are cocartesian. Next, we consider the
ring map
$$
A' = A'_1 \times_A A'_2 \longrightarrow
B' =
(A'_1 \otimes_{A_1} B_1)_{h_1} \times_B
(A'_2 \otimes_{A_1} B_2)_{h_2}
$$
By
More on Algebra, Lemma \href{https://stacks.math.columbia.edu/tag/07RU}{lemma module over fibre product (Tag 07RU)}
we have $A'_1 \otimes_{A'} B' = (A'_1 \otimes_{A_1} B_1)_{h_1}$
and $A'_2 \otimes_{A'} B' = (A'_2 \otimes_{A_2} B_2)_{h_2}$. In particular
the fibres of the morphism $Y' = \Spec(B') \to \Spec(A') = X'$ are
open subschemes of base changes of the fibres of the maps $Y_i \to X_i$, hence
finite and local complete intersections (see
Discriminants, Lemma \href{https://stacks.math.columbia.edu/tag/0BWE}{lemma syntomic quasi finite (Tag 0BWE)}).
By More on Algebra, Lemma
\href{https://stacks.math.columbia.edu/tag/08KQ}{lemma properties algebras over fibre product (Tag 08KQ)}
the ring map $A' \to B'$ is flat and of finite presentation.
Thus by Discriminants, Lemma
\href{https://stacks.math.columbia.edu/tag/0BWE}{lemma syntomic quasi finite (Tag 0BWE)} part (3)
we conclude that $Y'/X'$ is an object of $\mathcal{C}$.
Consider now the commutative diagram
$$
\xymatrix{
& & Y/X \ar[ld] \ar@/_2pc/[lldd]_{b_1/a_1} \ar[rd] \ar@/^2pc/[rrdd]^{b_2/a_2} \\
& Y'_1/X'_1 \ar[rd] \ar[ld] & & Y'_2/X'_2 \ar[ld] \ar[rd] \\
Y_1/X_1 & & Y'/X' & & Y_2/X_2
}
$$
where $Y'_i/X'_i$ corresponds to $A'_i \to (A'_i \otimes_{A_i} B_i)_{h_i}$.
Note that $Y'_i/X'_i$ is an object of $\mathcal{C}_{nice}$ as it
obtained by taking an open subscheme of the product of an
infinite dimensional affine space with $Y_i/X_i$; small detail omitted.
In particular, the pullback of $c^p_{Y_i/X_i}$ via $(b_i/a_i)$ is the same
as the pullback of $c^p_{Y'_i/X'_i}$ via $Y/X \to Y'_i/X'_i$.
Now we would be done if $Y'/X'$ is an object of $\mathcal{C}_{nice}$,
but this is almost never the case. Namely, then pulling back $c^p_{Y'/X'}$
around the two sides of the square, we would obtain the desired conclusion.
To get around the problem that $Y'/X'$ is not in $\mathcal{C}_{nice}$
we note the arguments above show that, after possibly shrinking all
of the schemes $X, Y, X'_1, Y'_1, X'_2, Y'_2, X', Y'$ we can find some
$n, d \geq 1$, and extend the diagram like so:
$$
\xymatrix{
& Y/X \ar[ld] \ar[rd] \\
Y'_1/X'_1 \ar[rd] & & Y'_2/X'_2 \ar[ld] \\
& Y'/X' \ar[d] \\
& Y_{n, d}/X_{n, d}
}
$$
and then we can use the already given argument by pulling
back from $c^p_{Y_{n, d}/X_{n, d}}$. This finishes the proof.
\end{proof}


\begin{proposition}
\label{discriminant-proposition-tate-map}
There exists a unique rule that to every locally quasi-finite syntomic
morphism of locally Noetherian schemes $Y \to X$ assigns an isomorphism
$$
c_{Y/X} : \det(\NL_{Y/X}) \longrightarrow \omega_{Y/X}
$$
satisfying the following two properties
\begin{enumerate}
\item the section $\delta(\NL_{Y/X})$ is mapped to $\tau_{Y/X}$, and
\item the rule is compatible with restriction to opens and with
base change.
\end{enumerate}
\end{proposition}

\begin{proof}
Let us reformulate the statement of the proposition. Consider the category
$\mathcal{C}$ whose objects, denoted $Y/X$, are locally quasi-finite syntomic
morphism $Y \to X$ of locally Noetherian schemes and whose morphisms
$b/a : Y'/X' \to Y/X$ are commutative diagrams
$$
\xymatrix{
Y' \ar[d] \ar[r]_b & Y \ar[d] \\
X' \ar[r]^a & X
}
$$
which induce an isomorphism of $Y'$ with an open subscheme of
$X' \times_X Y$. The proposition means that for every object
$Y/X$ of $\mathcal{C}$ we have an isomorphism
$c_{Y/X} : \det(\NL_{Y/X}) \to \omega_{Y/X}$
with $c_{Y/X}(\delta(\NL_{Y/X})) = \tau_{Y/X}$
and for every morphism $b/a : Y'/X' \to Y/X$ of $\mathcal{C}$ we have
$b^*c_{Y/X} = c_{Y'/X'}$ via the identifications
$b^*\det(\NL_{Y/X}) = \det(\NL_{Y'/X'})$ and
$b^*\omega_{Y/X} = \omega_{Y'/X'}$ described above.

\medskip\noindent
Given $Y/X$ in $\mathcal{C}$ and $y \in Y$ we can find
an affine opens $V \subset Y$ and $U \subset X$ with
$y \in V$ and $f(V) \subset U$
such that there exists some isomorphism
$$
\det(\NL_{Y/X})|_V \longrightarrow \omega_{Y/X}|_V
$$
mapping $\delta(\NL_{Y/X})|_V$ to $\tau_{Y/X}|_V$. This follows
from picking affine opens as in
Lemma \href{https://stacks.math.columbia.edu/tag/0BWE}{lemma syntomic quasi finite (Tag 0BWE)} part (5), the affine
local description of $\delta(\NL_{Y/X})$ in
Remark \href{https://stacks.math.columbia.edu/tag/0FKC}{remark local description delta (Tag 0FKC)}, and
Lemma \href{https://stacks.math.columbia.edu/tag/0BWH}{lemma different quasi finite complete intersection (Tag 0BWH)}.
If the annihilator of the section $\tau_{Y/X}$ is zero, then
these local maps are unique and automatically glue. Hence if the annihilator
of $\tau_{Y/X}$ is zero, then there is a unique isomorphism
$c_{Y/X} : \det(\NL_{Y/X}) \to \omega_{Y/X}$ with
$c_{Y/X}(\delta(\NL_{Y/X})) = \tau_{Y/X}$.
If $b/a : Y'/X' \to Y/X$ is a morphism of $\mathcal{C}$
and the annihilator of $\tau_{Y'/X'}$ is zero as well,
then $b^*c_{Y/X}$ is the unique isomorphism
$c_{Y'/X'} : \det(\NL_{Y'/X'}) \to \omega_{Y'/X'}$ with
$c_{Y'/X'}(\delta(\NL_{Y'/X'})) = \tau_{Y'/X'}$.
This follows formally from the fact that
$b^*\delta(\NL_{Y/X}) = \delta(\NL_{Y'/X'})$ and
$b^*\tau_{Y/X} = \tau_{Y'/X'}$.

\medskip\noindent
We can summarize the results of the previous paragraph as follows.
Let $\mathcal{C}_{nice} \subset \mathcal{C}$ denote the
full subcategory of $Y/X$ such that the annihilator of
$\tau_{Y/X}$ is zero. Then we have solved the problem
on $\mathcal{C}_{nice}$. For $Y/X$ in $\mathcal{C}_{nice}$
we continue to denote $c_{Y/X}$ the solution we've just found.

\medskip\noindent
Consider morphisms
$$
Y_1/X_1 \xleftarrow{b_1/a_1} Y/X \xrightarrow{b_2/a_2} Y_2/X_2
$$
in $\mathcal{C}$ such that $Y_1/X_1$ and $Y_2/X_2$ are objects
of $\mathcal{C}_{nice}$. {\bf Claim.} $b_1^*c_{Y_1/X_1} = b_2^*c_{Y_2/X_2}$.
We will first show that the claim implies the proposition
and then we will prove the claim.

\medskip\noindent
Let $d, n \geq 1$ and consider the locally
quasi-finite syntomic morphism $Y_{n, d} \to X_{n, d}$
constructed in Example \href{https://stacks.math.columbia.edu/tag/0FK8}{example universal quasi finite syntomic (Tag 0FK8)}.
Then $Y_{n, d}$ is an irreducible regular scheme and the
morphism $Y_{n, d} \to X_{n, d}$ is locally quasi-finite syntomic
and \'etale over a dense open, see
Lemma \href{https://stacks.math.columbia.edu/tag/0FK9}{lemma universal quasi finite syntomic etale (Tag 0FK9)}.
Thus $\tau_{Y_{n, d}/X_{n, d}}$ is nonzero for example by
Lemma \href{https://stacks.math.columbia.edu/tag/0BW9}{lemma different ramification (Tag 0BW9)}. Now a nonzero section
of an invertible module over an irreducible regular scheme
has vanishing annihilator. Thus
$Y_{n, d}/X_{n, d}$ is an object of $\mathcal{C}_{nice}$.

\medskip\noindent
Let $Y/X$ be an arbitrary object of $\mathcal{C}$. Let $y \in Y$.
By Lemma \href{https://stacks.math.columbia.edu/tag/0FKA}{lemma locally comes from universal (Tag 0FKA)} we can find
$n, d \geq 1$ and morphisms
$$
Y/X \leftarrow V/U \xrightarrow{b/a} Y_{n, d}/X_{n, d}
$$
of $\mathcal{C}$ such that $V \subset Y$ and $U \subset X$ are open.
Thus we can pullback the canonical morphism $c_{Y_{n, d}/X_{n, d}}$
constructed above by $b$ to $V$. The claim guarantees these local
isomorphisms glue! Thus we get a well defined global isomorphism
$c_{Y/X} : \det(\NL_{Y/X}) \to \omega_{Y/X}$ with
$c_{Y/X}(\delta(\NL_{Y/X})) = \tau_{Y/X}$.
If $b/a : Y'/X' \to Y/X$ is a morphism of $\mathcal{C}$, then
the claim also implies that the similarly constructed map
$c_{Y'/X'}$ is the pullback by $b$ of the locally constructed
map $c_{Y/X}$. Thus it remains to prove the claim.

\medskip\noindent
In the rest of the proof we prove the claim. We may pick a point
$y \in Y$ and prove the maps agree in an open neighbourhood of $y$.
Thus we may replace $Y_1$, $Y_2$ by open neighbourhoods of the
image of $y$ in $Y_1$ and $Y_2$. Therefore we may assume
$Y, X, Y_1, X_1, Y_2, X_2$ are affine, say they are the spectra
of rings $B, A, B_1, A_1, B_2, A_2$. Picture
$$
\xymatrix{
B_1 \ar[r] & B & B_2 \ar[l] \\
A_1 \ar[u] \ar[r] & A \ar[u] & A_2 \ar[l] \ar[u]
}
$$
By assumption the spectrum of $B$ is an affine open of
both the spectrum of $A \otimes_{A_1} B_1$ and $A \otimes_{A_2} B_2$.
Shrinking more we may assume there exist elements
$g_i \in A \otimes_{A_i} B_i$ such that our maps give isomorphisms
$(A \otimes_{A_i} B_i)_{g_i} = B$, see
Properties, Lemma \href{https://stacks.math.columbia.edu/tag/0H9B}{lemma standard open two affines (Tag 0H9B)}.
Let $x_\alpha$, $y_\beta$ be a sufficiently large collection of
variables such that we may choose surjections
$$
A'_1 = A_1[x_\alpha] \to A
\quad\text{and}\quad
A'_2 = A_2[y_\beta] \to A
$$
of $A_1$ and $A_2$-algebras. Then we can choose lifts
$h_i \in A'_i \otimes_{A_i} B_i$ of $g_i \in A \otimes_{A_i} B_i$
and we consider the diagram
$$
\xymatrix{
(A'_1 \otimes_{A_1} B_1)_{h_1} \ar[r] & B &
(A'_2 \otimes_{A_1} B_2)_{h_2} \ar[l] \\
A'_1 \ar[u] \ar[r] & A \ar[u] & A'_2 \ar[l] \ar[u]
}
$$
By construction the two squares are cocartesian. Next, we consider the
ring map
$$
A' = A'_1 \times_A A'_2 \longrightarrow
B' =
(A'_1 \otimes_{A_1} B_1)_{h_1} \times_B
(A'_2 \otimes_{A_1} B_2)_{h_2}
$$
By
More on Algebra, Lemma \href{https://stacks.math.columbia.edu/tag/07RU}{lemma module over fibre product (Tag 07RU)}
we have $A'_1 \otimes_{A'} B' = (A'_1 \otimes_{A_1} B_1)_{h_1}$
and $A'_2 \otimes_{A'} B' = (A'_2 \otimes_{A_2} B_2)_{h_2}$. In particular
the fibres of the morphism $Y' = \Spec(B') \to \Spec(A') = X'$ are
open subschemes of base changes of the fibres of the maps $Y_i \to X_i$,
hence finite and local complete intersections
(see Lemma \href{https://stacks.math.columbia.edu/tag/0BWE}{lemma syntomic quasi finite (Tag 0BWE)}).
By More on Algebra, Lemma
\href{https://stacks.math.columbia.edu/tag/08KQ}{lemma properties algebras over fibre product (Tag 08KQ)}
the ring map $A' \to B'$ is flat and of finite presentation.
Thus by Lemma \href{https://stacks.math.columbia.edu/tag/0BWE}{lemma syntomic quasi finite (Tag 0BWE)} part (3)
we conclude that $Y' \to X'$ is locally quasi-finite and syntomic.
Thus we obtain the commutative diagram
$$
\xymatrix{
& & Y/X \ar[ld] \ar@/_2pc/[lldd]_{b_1/a_1} \ar[rd] \ar@/^2pc/[rrdd]^{b_2/a_2} \\
& Y'_1/X'_1 \ar[rd] \ar[ld] & & Y'_2/X'_2 \ar[ld] \ar[rd] \\
Y_1/X_1 & & Y'/X' & & Y_2/X_2
}
$$
where $Y'_i/X'_i$ corresponds to $A'_i \to (A'_i \otimes_{A_i} B_i)_{h_i}$.

\medskip\noindent
In this paragraph, we fix the issue that $A'$ and $A'_i$ may not be Noetherian
by a limit argument; we suggest the reader skip this.
We can find a finitely generated $\mathbf{Z}$-subalgebra
$A'' \subset A'$ and a quasi-finite and syntomic ring map
$A'' \to B''$ such that $B' = A' \otimes_{A''} B''$.
This follows from Limits, Lemmas
\href{https://stacks.math.columbia.edu/tag/01ZM}{lemma descend finite presentation (Tag 01ZM)},
\href{https://stacks.math.columbia.edu/tag/0C3L}{lemma descend syntomic (Tag 0C3L)}, and
\href{https://stacks.math.columbia.edu/tag/094M}{lemma descend quasi finite (Tag 094M)}.
Then the ring maps $A'' \to A'_1 = A_1[x_\alpha]$ will factor through a finite
polynomial subalgebra $A_i[x_1, \ldots, x_n$ and the ring map
$A'' \to A'_2 = A_2[y_\beta]$ will factor through $A_2[y_1, \ldots, y_m]$.
After increasing the sets of variables, we may also assume that
$h_1 \in A_1[x_1, \ldots, x_n] \otimes_{A_1} B_1$ and
$h_2 \in A_2[y_1, \ldots, y_m] \otimes_{A_2} B_2$.
Increasing the sets of variables one more time, we can assume
that the maps $B'' \to (A'_i \otimes_{A_i} B_i)_{h_i}$
factor through $(A_1[x_1, \ldots, x_n] \otimes_{A_1} B_1)_{h_1}$
and $(A_2[y_1, \ldots, y_m] \otimes_{A_2} B_2)_{h_2}$.
Replacing $Y'$, $X'$, $Y'_1$, $X'_1$, $Y'_2$, $X'_2$ by the spectra of
$B''$, $A''$, $(A_1[x_1, \ldots, x_n] \otimes_{A_1} B_1)_{h_1}$,
$A_1[x_1, \ldots, x_n]$, $(A_2[y_1, \ldots, y_m] \otimes_{A_2} B_2)_{h_2}$,
$A_2[y_1, \ldots, y_m]$, we see that we get a diagram as above
where all schemes are Noetherian.

\medskip\noindent
Note that $Y'_i/X'_i$ is an object of $\mathcal{C}_{nice}$ as it
obtained by taking an open subscheme of the product of an affine space
with $Y_i/X_i$. In particular, the pullback of
$c_{Y_i/X_i}$ via $(b_i/a_i)$ is the same as the pullback of
$c_{Y'_i/X'_i}$ via $Y/X \to Y'_i/X'_i$.
Now we would be done if $Y'/X'$ is an object of $\mathcal{C}_{nice}$
(but this is likely not the case). Namely, then pulling back $c_{Y'/X'}$
around the two sides of the square, we would obtain the desired conclusion.
To get around the problem that $Y'/X'$ is not in $\mathcal{C}_{nice}$
we note the arguments above show that, after possibly shrinking all
of the schemes $X, Y, X'_1, Y'_1, X'_2, Y'_2, X', Y'$ we can find some
$n, d \geq 1$, and extend the diagram like so:
$$
\xymatrix{
& Y/X \ar[ld] \ar[rd] \\
Y'_1/X'_1 \ar[rd] & & Y'_2/X'_2 \ar[ld] \\
& Y'/X' \ar[d] \\
& Y_{n, d}/X_{n, d}
}
$$
and then we can use the already given argument by pulling
back from $c_{Y_{n, d}/X_{n, d}}$. This finishes the proof.
\end{proof}


\section{Finite morphisms}
\label{discriminant-section-finite-morphisms}

\noindent
In this section we collect some observations about the
constructions in the previous sections for finite morphisms.
Let $f : Y \to X$ be a finite morphism of locally Noetherian schemes.
Let $\omega_{Y/X}$ be as in
Remark \href{https://stacks.math.columbia.edu/tag/0BVG}{remark relative dualizing for quasi finite (Tag 0BVG)}.

\medskip\noindent
The first remark is that
$$
f_*\omega_{Y/X} = \SheafHom_{\mathcal{O}_X}(f_*\mathcal{O}_Y, \mathcal{O}_X)
$$
as sheaves of $f_*\mathcal{O}_Y$-modules. Since $f$ is affine, this
formula uniquely characterizes $\omega_{Y/X}$, see
Morphisms, Lemma \href{https://stacks.math.columbia.edu/tag/01SB}{lemma affine equivalence modules (Tag 01SB)}.
The formula holds because for $\Spec(A) = U \subset X$ affine open, the
inverse image $V = f^{-1}(U)$ is the spectrum of a finite $A$-algebra
$B$ and hence
$$
H^0(U, f_*\omega_{Y/X}) =
H^0(V, \omega_{Y/X}) =
\omega_{B/A} =
\Hom_A(B, A) =
H^0(U, \SheafHom_{\mathcal{O}_X}(f_*\mathcal{O}_Y, \mathcal{O}_X))
$$
by construction. In particular, we obtain a canonical evaluation map
$$
f_*\omega_{Y/X} \longrightarrow \mathcal{O}_X
$$
which is given by evaluation at $1$ if we think of $f_*\omega_{Y/X}$
as the sheaf $\SheafHom_{\mathcal{O}_X}(f_*\mathcal{O}_Y, \mathcal{O}_X)$.

\medskip\noindent
The second remark is that using the evaluation map we obtain
canonical identifications
$$
\Hom_Y(\mathcal{F}, f^*\mathcal{G} \otimes_{\mathcal{O}_Y} \omega_{Y/X})
=
\Hom_X(f_*\mathcal{F}, \mathcal{G})
$$
functorially in the quasi-coherent module $\mathcal{F}$ on $Y$
and the finite locally free module $\mathcal{G}$ on $X$.
If $\mathcal{G} = \mathcal{O}_X$ this follows immediately
from the above and
Algebra, Lemma \href{https://stacks.math.columbia.edu/tag/08YP}{lemma adjoint hom restrict (Tag 08YP)}.
For general $\mathcal{G}$ we can use the same lemma and the
isomorphisms
$$
f_*(f^*\mathcal{G} \otimes_{\mathcal{O}_Y} \omega_{Y/X}) =
\mathcal{G} \otimes_{\mathcal{O}_X}
\SheafHom_{\mathcal{O}_X}(f_*\mathcal{O}_Y, \mathcal{O}_X) =
\SheafHom_{\mathcal{O}_X}(f_*\mathcal{O}_Y, \mathcal{G})
$$
of $f_*\mathcal{O}_Y$-modules where the first equality is the
projection formula
(Cohomology, Lemma \href{https://stacks.math.columbia.edu/tag/01E8}{lemma projection formula (Tag 01E8)}).
An alternative is to prove the formula affine locally by
direct computation.

\medskip\noindent
The third remark is that if $f$ is in addition flat, then the
composition
$$
f_*\mathcal{O}_Y \xrightarrow{f_*\tau_{Y/X}} f_*\omega_{Y/X}
\longrightarrow \mathcal{O}_X
$$
is equal to the trace map $\text{Trace}_f$ discussed in
Section \href{https://stacks.math.columbia.edu/tag/0BVH}{section discriminant (Tag 0BVH)}. This follows immediately by
looking over affine opens.

\medskip\noindent
The fourth remark is that if $f$ is flat and $X$ Noetherian, then
we obtain
$$
\Hom_Y(K, Lf^*M \otimes_{\mathcal{O}_Y} \omega_{Y/X})
=
\Hom_X(Rf_*K, M)
$$
for any $K$ in $D_\QCoh(\mathcal{O}_Y)$ and $M$ in $D_\QCoh(\mathcal{O}_X)$.
This follows from the material in
Duality for Schemes, Section \href{https://stacks.math.columbia.edu/tag/0E4H}{section proper flat (Tag 0E4H)},
but can be proven directly in this case as follows.
First, if $X$ is affine, then it holds by
Dualizing Complexes, Lemmas \href{https://stacks.math.columbia.edu/tag/0A70}{lemma right adjoint (Tag 0A70)} and
\href{https://stacks.math.columbia.edu/tag/0BZE}{lemma RHom is tensor special (Tag 0BZE)}\footnote{There is a
simpler proof of this lemma in our case.} and
Derived Categories of Schemes, Lemma \href{https://stacks.math.columbia.edu/tag/06Z0}{lemma affine compare bounded (Tag 06Z0)}.
Then we can use the induction principle
(Cohomology of Schemes, Lemma \href{https://stacks.math.columbia.edu/tag/08DR}{lemma induction principle (Tag 08DR)})
and Mayer-Vietoris
(in the form of Cohomology, Lemma \href{https://stacks.math.columbia.edu/tag/08BW}{lemma mayer vietoris hom (Tag 08BW)})
to finish the proof.









\section*{Primary problem: Tag 0FLB}
\begin{lemma}
\label{lemma-Garel}
There exists a unique rule that to every finite syntomic
morphism of schemes $f : Y \to X$ assigns $\mathcal{O}_X$-module maps
$$
\Theta^p_{Y/X} :
f_*\Omega^p_{Y/\mathbf{Z}}
\longrightarrow
\Omega^p_{X/\mathbf{Z}}
$$
satisfying the following properties
\begin{enumerate}
\item the composition with
$\Omega^p_{X/\mathbf{Z}} \otimes_{\mathcal{O}_X} f_*\mathcal{O}_Y
\to f_*\Omega^p_{Y/\mathbf{Z}}$ is equal to
$\text{id} \otimes \text{Trace}_f$
where $\text{Trace}_f : f_*\mathcal{O}_Y \to \mathcal{O}_X$
is the map from
Discriminants, Section \href{https://stacks.math.columbia.edu/tag/0BVH}{section discriminant (Tag 0BVH)},
\item the rule is compatible with base change.
\end{enumerate}
\end{lemma}

\begin{proof}
First, assume that $X$ is locally Noetherian. By
Lemma \ref{lemma-Garel-upstairs} we have a canonical map
$$
c^p_{Y/X} : \Omega_{Y/\mathbf{Z}}^p
\longrightarrow
f^*\Omega_{X/\mathbf{Z}}^p \otimes_{\mathcal{O}_Y} \det(\NL_{Y/X})
$$
By Discriminants, Proposition \ref{discriminant-proposition-tate-map}
we have a canonical isomorphism
$$
c_{Y/X} : \det(\NL_{Y/X}) \to \omega_{Y/X}
$$
mapping $\delta(\NL_{Y/X})$ to $\tau_{Y/X}$. Combined these maps give
$$
c^p_{Y/X} \otimes c_{Y/X} :
\Omega_{Y/\mathbf{Z}}^p
\longrightarrow
f^*\Omega_{X/\mathbf{Z}}^p \otimes_{\mathcal{O}_Y} \omega_{Y/X}
$$
By Discriminants, Section \ref{discriminant-section-finite-morphisms}
this is the same thing as a map
$$
\Theta_{Y/X}^p :
f_*\Omega_{Y/\mathbf{Z}}^p
\longrightarrow
\Omega_{X/\mathbf{Z}}^p
$$
Recall that the relationship between $c^p_{Y/X} \otimes c_{Y/X}$
and $\Theta_{Y/X}^p$ uses the evaluation map
$f_*\omega_{Y/X} \to \mathcal{O}_X$
which sends $\tau_{Y/X}$ to $\text{Trace}_f(1)$, see
Discriminants, Section \ref{discriminant-section-finite-morphisms}.
Hence property (1) holds. Property (2) holds for base changes by
$X' \to X$ with $X'$ locally Noetherian because both $c^p_{Y/X}$ and
$c_{Y/X}$ are compatible with such base changes. For $f : Y \to X$
finite syntomic and $X$ locally Noetherian,
we will continue to denote $\Theta^p_{Y/X}$ the solution we've just found.

\medskip\noindent
Uniqueness. Suppose that we have a finite syntomic morphism
$f: Y \to X$ such that $X$ is smooth over $\Spec(\mathbf{Z})$
and $f$ is \'etale over a dense open of $X$. We claim that
in this case $\Theta^p_{Y/X}$ is uniquely determined by property (1).
Namely, consider the maps
$$
\Omega^p_{X/\mathbf{Z}} \otimes_{\mathcal{O}_X} f_*\mathcal{O}_Y \to
f_*\Omega^p_{Y/\mathbf{Z}} \to
\Omega^p_{X/\mathbf{Z}}
$$
The sheaf $\Omega^p_{X/\mathbf{Z}}$ is torsion free (by the assumed
smoothness), hence it suffices to check that the restriction of
$\Theta^p_{Y/X}$ is uniquely determined over the dense open over
which $f$ is \'etale, i.e., we may assume $f$ is \'etale.
However, if $f$ is \'etale, then
$f^*\Omega_{X/\mathbf{Z}} = \Omega_{Y/\mathbf{Z}}$
hence the first arrow in the displayed equation is an isomorphism.
Since we've pinned down the composition, this guarantees uniqueness.

\medskip\noindent
Let $f : Y \to X$ be a finite syntomic morphism of locally Noetherian schemes.
Let $x \in X$. By Discriminants, Lemma
\href{https://stacks.math.columbia.edu/tag/0FL4}{lemma locally comes from universal finite (Tag 0FL4)}
we can find $d \geq 1$ and a commutative diagram
$$
\xymatrix{
Y \ar[d] &
V \ar[d] \ar[l] \ar[r] &
V_d \ar[d] \\
X &
U \ar[l] \ar[r] &
U_d
}
$$
such that $x \in U \subset X$ is open, $V = f^{-1}(U)$
and $V = U \times_{U_d} V_d$. Thus $\Theta^p_{Y/X}|_V$
is the pullback of the map $\Theta^p_{V_d/U_d}$.
However, by the discussion on uniqueness above and
Discriminants, Lemmas
\href{https://stacks.math.columbia.edu/tag/0FL1}{lemma universal finite syntomic smooth (Tag 0FL1)} and
\href{https://stacks.math.columbia.edu/tag/0FL2}{lemma universal finite syntomic etale (Tag 0FL2)}
the map $\Theta^p_{V_d/U_d}$ is uniquely determined
by the requirement (1). Hence uniqueness holds.

\medskip\noindent
At this point we know that we have existence and uniqueness
for all finite syntomic morphisms $Y \to X$ with $X$ locally Noetherian.
We could now give an argument similar to the proof of
Lemma \ref{lemma-Garel-upstairs} to extend to general $X$.
However, instead it is possible to directly use
absolute Noetherian approximation to finish the proof.
Namely, to construct $\Theta^p_{Y/X}$
it suffices to do so Zariski locally on $X$ (provided we also
show the uniqueness). Hence we may assume $X$ is affine (small
detail omitted). Then we can write $X = \lim_{i \in I} X_i$
as the limit over a directed set $I$ of Noetherian affine schemes.
By Algebra, Lemma \href{https://stacks.math.columbia.edu/tag/05N9}{lemma colimit category fp algebras (Tag 05N9)}
we can find $0 \in I$ and a finitely
presented morphism of affines $f_0 : Y_0 \to X_0$ whose base change to
$X$ is $Y \to X$. After increasing $0$ we may assume $Y_0 \to X_0$
is finite and syntomic, see
Algebra, Lemma \href{https://stacks.math.columbia.edu/tag/0C33}{lemma colimit lci (Tag 0C33)} and
\href{https://stacks.math.columbia.edu/tag/07RG}{lemma colimit finite (Tag 07RG)}. For $i \geq 0$ also the
base change $f_i : Y_i = Y_0 \times_{X_0} X_i \to X_i$ is finite syntomic.
Then
$$
\Gamma(X, f_*\Omega^p_{Y/\mathbf{Z}}) =
\Gamma(Y, \Omega^p_{Y/\mathbf{Z}}) =
\colim_{i \geq 0} \Gamma(Y_i, \Omega^p_{Y_i/\mathbf{Z}}) =
\colim_{i \geq 0} \Gamma(X_i, f_{i, *}\Omega^p_{Y_i/\mathbf{Z}})
$$
Hence we can (and are forced to) define $\Theta^p_{Y/X}$ as the colimit
of the maps $\Theta^p_{Y_i/X_i}$. This map is compatible with any
cartesian diagram
$$
\xymatrix{
Y' \ar[r] \ar[d] & Y \ar[d] \\
X' \ar[r] & X
}
$$
with $X'$ affine as we know this for the case of Noetherian affine schemes
by the arguments given above (small detail omitted; hint: if we also
write $X' = \lim_{j \in J} X'_j$ then for every $i \in I$ there is a $j \in J$
and a morphism $X'_j \to X_i$ compatible with the morphism $X' \to X$).
This finishes the proof.
\end{proof}
\end{document}
